\chapter{动量映射与辛商}

% [原文 p.58]
设辛流形 $(M,\omega)$ 与其上的 Hamilton 函数 $H$ 具有关于 Lie 群 $G$ 作用的对称性，则自然会想到用轨道空间 $M/G$ 替换 $M$，以减少系统的自由度。然而一般地 $M/G$ 并不是辛流形。若 $G$ 是奇数维的，则 $M/G$ 也是奇数维的，本来就无法具有辛结构。

辛几何中正确的商的概念，是通过注意到下面这一点而得到的：Hamilton 系统若具有对称性，则由 Noether 定理，相应的动量守恒。也就是说，先取动量取定值的部分空间，再用该部分空间上的群作用取商，如此得到的商称为辛商。

为把这一点精确地表述出来，需要 Lie 群对辛流形的 Hamilton 作用及其动量映射的概念。在这一过程中，4.4 节所见 Lie 代数的对偶空间上的 Poisson 结构会自然出现。

辛商与代数几何中的几何不变量理论商（geometric invariant theory quotient）密切相关。定义辛商，不仅需要辛流形 $(M,\omega)$ 及其上的 Lie 群 $G$ 的 Hamilton 作用，还必须选取 $\mathfrak{g}^*$ 的一个元素 $\xi$。这与 GIT 商中直线丛的 G 线性化（$G$-linearization）的选取相对应。改变 $\xi$ 所引起的辛商的变化，在代数几何中称为 VGIT（variation of GIT quotient），在理论上与应用上都很重要。VGIT 的典型例子可在第 8 章所讨论的环面几何中见到。
% [Tutor Note: 背景——把辛商与 GIT 商对应起来的是 Kempf--Ness 定理；改变 $\xi$ 即改变 GIT 中直线丛的线性化，其变化（VGIT）是 Thaddeus 等人研究的主题。]

\section{Hamilton 作用}

设 Lie 群 $G$ 作用于辛流形 $(M,\omega)$；
% [原文 p.59]
\begin{equation}
  a\colon G\times M\to M,\qquad (g,x)\mapsto g\cdot x.
  \tag{5.1.1}
\end{equation}

对作用 $a$ 微分，可得映射 $\mathfrak{g}\to\Gamma(TM)$；把由此得到的 $Y\in\mathfrak{g}$ 的像记作 $\underline{Y}\in\Gamma(TM)$，称为 $Y$ 的\emph{基本向量场}（fundamental vector field）。假定对任意 $Y\in\mathfrak{g}$，存在函数 $\mu_Y\colon M\to\R$，使 $Y$ 的基本向量场 $\underline{Y}$ 成为以 $\mu_Y$ 为 Hamilton 函数的 Hamilton 向量场；
\begin{equation}
  X_{\mu_Y}=\underline{Y}.
  \tag{5.1.2}
\end{equation}
由 Hamilton 向量场的定义，(5.1.2) 可改写为
\begin{equation}
  d\mu_Y=\iota_{\underline{Y}}\omega.
  \tag{5.1.3}
\end{equation}

\takeaway{动量映射把 Lie 代数的每个元素对应到基本向量场的 Hamilton 函数，从而把连续对称性编码为守恒量。}
% [Tutor Note: 难点——(5.1.2) 之所以能对每个 $Y$ 要求 $\underline{Y}$ 是 Hamilton 向量场，是因为 $G$ 保 $\omega$，故 $\iota_{\underline{Y}}\omega$ 是闭形式；在 $M$ 单连通时才自动恰当。]

对给定的向量场，以它为 Hamilton 向量场的函数在相差一个常数的意义下唯一确定。对 $\mathfrak{g}$ 的每个元素 $Y$ 各选定一个这样的常数并固定，便得到映射
\begin{equation}
  \mathfrak{g}\to C^\infty(M),\qquad Y\mapsto\mu_Y.
  \tag{5.1.4}
\end{equation}
进一步，适当选取常数，可使 (5.1.4) 成为线性映射。此时，存在 $M$ 上的 $\mathfrak{g}^*$ 值 $C^\infty$ 函数
\begin{equation}
  \mu\colon M\to\mathfrak{g}^*,
  \tag{5.1.5}
\end{equation}
使对任意 $Y\in\mathfrak{g}$ 有
\begin{equation}
  \mu_Y=\langle\mu,Y\rangle
  \tag{5.1.6}
\end{equation}
成立。此映射由命题 2.8.2 满足
\begin{align}
  X_{\{\mu_Y,\mu_Z\}}&=[X_{\mu_Y},X_{\mu_Z}] \tag{5.1.7}\\
  &=[Y,Z] \tag{5.1.8}\\
  &=X_{\mu_{[Y,Z]}}, \tag{5.1.9}
\end{align}
但以给定向量场为 Hamilton 向量场的 Hamilton 函数有加上常数的自由度，故对任意 $Y,Z\in\mathfrak{g}$ 未必满足
\begin{equation}
  \mu_{[Y,Z]}=\{\mu_Y,\mu_Z\},\qquad \forall Y,Z\in\mathfrak{g}
  \tag{5.1.10}
\end{equation}
% [Tutor Note: 校勘——(5.1.8) 的右端原书未加下划线，按 (5.1.2) 应理解为 $[\underline{Y},\underline{Z}]$；此处照录原文。]
成立。

\begin{definition}{5.1.1}
  辛流形 $(M,\omega)$ 上的连通 Lie 群 $G$ 的 \textbf{Hamilton 作用}（Hamiltonian action），是指 $G$ 对 $M$ 的作用与 $M$ 上 $C^\infty$ 函数
  \begin{equation}
    \mu\colon M\to\mathfrak{g}^*
    \tag{5.1.11}
  \end{equation}
  的组合，且满足以下性质：
  % [原文 p.60]
  \begin{enumerate}
    \item $G$ 对 $M$ 的作用保持辛形式 $\omega$，即对任意 $g\in G$ 有 $g^*\omega=\omega$ 成立。
    \item 对任意 $Y\in\mathfrak{g}$，$Y$ 所生成的 $M$ 上的向量场 $\underline{Y}$ 是以 $\langle\mu,Y\rangle\colon M\to\R$ 为 Hamilton 函数的 Hamilton 向量场；
    \begin{equation}
      \iota_{\underline{Y}}\omega=d\langle\mu,Y\rangle.
      \tag{5.1.12}
    \end{equation}
    \item 映射 $\mu\colon M\to\mathfrak{g}^*$ 是 Poisson 映射，即对任意 $Y,Z\in\mathfrak{g}$ 有
    \begin{equation}
      \mu_{[Y,Z]}=\{\mu_Y,\mu_Z\}
      \tag{5.1.13}
    \end{equation}
    成立。
  \end{enumerate}
  此时，称 $\mu$ 为\emph{动量映射}（moment map）。
\end{definition}
% [Tutor Note: 难点——在 $\mathfrak{g}$ 的一组基上任意取常数再线性延拓，(5.1.4) 便成为线性映射；而 (5.1.13) 是额外要求，正是它定义了 Hamilton 作用而非一般的辛作用。]

\begin{proposition*}{5.1.2}
  动量映射的条件 (5.1.13) 与 $\mu$ 的等变性
  \begin{equation}
    \mu(g\cdot x)=\operatorname{Ad}^*_g(\mu(x)),\qquad \forall g\in G,\ x\in M
    \tag{5.1.14}
  \end{equation}
  等价。
\end{proposition*}

\begin{proof}
  对 $Y\in\mathfrak{g}$ 令 $g=\exp(tY)$，取 (5.1.14) 左端与 $Z\in\mathfrak{g}$ 的配对并在 $t=0$ 处微分，则由基本向量场的定义、(5.1.2) 以及 Poisson 括号的定义得
  \begin{align}
    \left.\frac{d}{dt}\mu_Z(e^{tY}\cdot x)\right|_{t=0}&=\underline{Y}\mu_Z(x) \tag{5.1.15}\\
    &=X_{\mu_Y}\mu_Z(x) \tag{5.1.16}\\
    &=\{\mu_Y,\mu_Z\}(x). \tag{5.1.17}
  \end{align}
  % [Tutor Note: 难点——(5.1.15) 中 $e^{tY}\cdot x$ 关于 $t$ 的导数正是基本向量场 $\underline{Y}$ 在 $x$ 处的值，故对 $\mu_Z$ 求导即 $\underline{Y}\mu_Z$。]
  另一方面，对 (5.1.14) 的右端在 $t=0$ 处微分得
  \begin{align}
    \left.\frac{d}{dt}\langle\operatorname{Ad}^*_{e^{tY}}\mu,Z\rangle\right|_{t=0}&=\langle\ad^*_Y\mu,Z\rangle \tag{5.1.18}\\
    &=\langle\mu,\ad_YZ\rangle \tag{5.1.19}\\
    &=\langle\mu,[Y,Z]\rangle \tag{5.1.20}\\
    &=\mu_{[Y,Z]}. \tag{5.1.21}
  \end{align}
  由于 $G$ 连通，(5.1.14) 在单位元 $g=e$ 处显然成立，于是由常微分方程解的唯一性可知：(5.1.14) 对任意 $g\in G$ 成立的充要条件是微分所得的 (5.1.13) 成立。
\end{proof}

下述 2 个命题由定义立即可得：

% [原文 p.61]
\begin{proposition}{5.1.3}
  设 Lie 群 $G$ 对 $(M,\omega)$ 的 Hamilton 作用已给出，其动量映射为 $\mu\colon M\to\mathfrak{g}^*$。此时，把 $G$ 的 Hamilton 作用限制到 $G$ 的任意闭子群 $H$ 上，便给出 $H$ 的 Hamilton 作用，其动量映射为 $\iota^*\circ\mu\colon M\to\mathfrak{h}^*$。其中，$\iota^*\colon\mathfrak{g}^*\to\mathfrak{h}^*$ 是 Lie 群的嵌入 $H\hookrightarrow G$ 所诱导的 Lie 代数嵌入 $\iota_*\colon\mathfrak{h}\hookrightarrow\mathfrak{g}$ 的对偶线性满射。
\end{proposition}

\begin{proposition}{5.1.4}
  设 Lie 群 $G_1$ 与 $G_2$ 分别对辛流形 $(M_1,\omega_1)$ 与 $(M_2,\omega_2)$ 有 Hamilton 作用，其动量映射分别由 $\mu_1\colon M_1\to\mathfrak{g}_1^*$ 与 $\mu_2\colon M_2\to\mathfrak{g}_2^*$ 给出。此时，$G_1\times G_2$ 对 $(M_1\times M_2,\ \pi_1^*\omega_1+\pi_2^*\omega_2)$ 有自然的 Hamilton 作用，其动量映射由
  \begin{equation}
    \mu_1+\mu_2\colon M_1\times M_2\to\mathfrak{g}_1^*\oplus\mathfrak{g}_2^*
    \tag{5.1.22}
  \end{equation}
  给出。
\end{proposition}

又，下述命题的证明也很容易：

\begin{proposition*}{5.1.5}
  对余伴随轨道 $\mathcal{O}\subset\mathfrak{g}^*$ 的余伴随作用 $G\curvearrowright\mathcal{O}$ 是 Hamilton 作用，其动量映射由嵌入映射
  \begin{equation}
    \iota\colon\mathcal{O}\hookrightarrow\mathfrak{g}^*
    \tag{5.1.23}
  \end{equation}
  给出。
\end{proposition*}

\begin{proof}
  只需对任意 $Y\in\mathfrak{g}$ 证明 $\mathcal{O}\subset\mathfrak{g}^*$ 上 1 次微分形式的等式
  \begin{equation}
    d\langle\iota,Y\rangle=\iota_{\underline{Y}}\omega.
    \tag{5.1.24}
  \end{equation}
  对任意 $\xi\in\mathcal{O}$，$Y$ 的基本向量场 $\underline{Y}$ 在 $\xi$ 处的值由
  \begin{equation}
    \underline{Y}_{\xi}=\ad^*_Y(\xi)\in\mathfrak{g}^*
    \tag{5.1.25}
  \end{equation}
  给出。又，任意 $\eta\in T_{\xi}\mathcal{O}$ 可由适当的 $Z\in\mathfrak{g}$ 表为
  \begin{equation}
    \eta=\ad^*_Z(\xi).
    \tag{5.1.26}
  \end{equation}
  取 (5.1.24) 的右端与 $\eta$ 的配对，得
  \begin{align}
    \omega(\underline{Y}_{\xi},\eta)&=\omega(\ad^*_Y(\xi),\ad^*_Z(\xi)) \tag{5.1.27}\\
    &=\langle\xi,[Y,Z]\rangle \tag{5.1.28}\\
    &=\langle\xi,-\ad_ZY\rangle \tag{5.1.29}\\
    &=\langle\ad^*_Z\xi,Y\rangle \tag{5.1.30}\\
    &=\langle\eta,Y\rangle, \tag{5.1.31}
  \end{align}
  而这显然等于 (5.1.24) 的左端与 $\eta$ 的配对。
\end{proof}
% [Tutor Note: 背景——$\mathcal{O}$ 上的辛形式即 Kostant--Kirillov--Souriau 形式，它使 $G$ 的余伴随轨道成为齐性辛流形，与 4.4 节 $\mathfrak{g}^*$ 上的 Lie--Poisson 结构相容。]

% [原文 p.62]
下述命题在构造 Hamilton 作用的例子时很重要：

\begin{proposition*}{5.1.6}
  恰当辛流形（exact symplectic manifold）$(M,\omega=d\theta)$ 上保持 $\theta$ 的 Lie 群作用都是 Hamilton 作用。
\end{proposition*}

\begin{proof}
  当 Lie 群 $G$ 保 $\theta$ 地作用于 $M$ 时，对任意 $Y\in\mathfrak{g}$ 有
  \begin{equation}
    \mathcal{L}_{\underline{Y}}\theta=0.
    \tag{5.1.32}
  \end{equation}
  令
  \begin{equation}
    \mu_Y=-\iota_{\underline{Y}}\theta,
    \tag{5.1.33}
  \end{equation}
  则它关于 $Y$ 是线性的，故存在某个 $\mu\colon M\to\mathfrak{g}^*$，使对任意 $Y\in\mathfrak{g}$ 有 $\mu_Y=\langle\mu,Y\rangle$ 成立。对任意 $Y,Z\in\mathfrak{g}$，
  \begin{align}
    d\mu_Y&=-d\iota_{\underline{Y}}\theta \tag{5.1.34}\\
    &=-(d\iota_{\underline{Y}}+\iota_{\underline{Y}}d)\theta+\iota_{\underline{Y}}d\theta \tag{5.1.35}\\
    &=-\mathcal{L}_{\underline{Y}}\theta+\iota_{\underline{Y}}d\theta \tag{5.1.36}\\
    &=\iota_{\underline{Y}}d\theta \tag{5.1.37}
  \end{align}
  且
  \begin{align}
    \mu_{[Y,Z]}&=-\iota_{\underline{[Y,Z]}}\theta \tag{5.1.38}\\
    &=-\iota_{[\underline{Y},\underline{Z}]}\theta \tag{5.1.39}\\
    &=-\iota_{\mathcal{L}_{\underline{Y}}\underline{Z}}\theta \tag{5.1.40}\\
    &=-\mathcal{L}_{\underline{Y}}(\iota_{\underline{Z}}\theta) \tag{5.1.41}\\
    &=\mathcal{L}_{\underline{Y}}(\mu_Z) \tag{5.1.42}\\
    &=\underline{Y}(\mu_Z) \tag{5.1.43}\\
    &=X_{\mu_Y}(\mu_Z) \tag{5.1.44}\\
    &=\{\mu_Y,\mu_Z\}, \tag{5.1.45}
  \end{align}
  故 $\mu$ 是 $G$ 作用的动量映射。
\end{proof}
% [Tutor Note: 难点——(5.1.34) 到 (5.1.37) 只是把 Cartan 同伦公式插入并整理；(5.1.40) 用到 $[\underline{Y},\underline{Z}]=\mathcal{L}_{\underline{Y}}\underline{Z}$，(5.1.41) 用到 $\mathcal{L}_{\underline{Y}}$ 是微分、与内乘的交换关系 $\mathcal{L}_{X}\iota_Y-\iota_Y\mathcal{L}_X=\iota_{[X,Y]}$。]

对余切丛的群作用，是命题 5.1.6 的典型应用例。

\begin{proposition*}{5.1.7}
  Lie 群 $G$ 对流形 $N$ 的作用，自然地提升为对余切丛 $(M,\omega)=(T^*N,d\theta)$ 的 Hamilton 作用。
\end{proposition*}

\begin{proof}
  由余切丛上 Liouville 形式 $\theta$ 的定义 (4.1.2) 可知，$\theta$ 关于 $G$ 作用不变，即对任意 $g\in G$ 有
  \begin{equation}
    g^*\theta=\theta
    \tag{5.1.46}
  \end{equation}
  成立，这一点是明显的。
\end{proof}

% [原文 p.63]
命题 5.1.6 的另一个典型应用例，是到辛向量空间的线性群作用。

\begin{proposition*}{5.1.8}
  辛向量空间 $\bigl(\C^n,\ \frac{\sqrt{-1}}{2}\sum_{i=1}^n\rd z_i\wedge\rd\bar z_i\bigr)$ 上酉群 $U(n)$ 的自然作用是 Hamilton 作用，其动量映射由
  \begin{equation}
    \mu\colon\C^n\to\mathfrak{u}(n),\qquad (z_i)_{i=1}^n\mapsto\left(\frac{\sqrt{-1}}{2}(z_i\bar z_j)\right)_{i,j=1}^n
    \tag{5.1.47}
  \end{equation}
  给出。
\end{proposition*}

\begin{proof}
  $\C^n$ 上的标准辛形式
  \begin{equation}
    \omega=\frac{\sqrt{-1}}{2}\sum_{i=1}^n\rd z_i\wedge\rd\bar z_i
    \tag{5.1.48}
  \end{equation}
  可由 $\C^n$ 上的 1 次微分形式
  \begin{equation}
    \theta=\frac{\sqrt{-1}}{4}\sum_{i=1}^n(z_i\wedge\rd\bar z_i-\bar z_i\wedge\rd z_i)
    \tag{5.1.49}
  \end{equation}
  表示为
  \begin{equation}
    \omega=d\theta.
    \tag{5.1.50}
  \end{equation}
  $U(n)$ 的元 $g=(g_{ij})_{i,j=1}^n$ 对 $\C^n$ 的作用由
  \begin{equation}
    z=(z_i)_{i=1}^n\mapsto g\cdot z=\left(\sum_{j=1}^n g_{ij}z_j\right)_{i=1}^n
    \tag{5.1.51}
  \end{equation}
  给出，故 $Y=(Y_{ij})_{i,j=1}^n\in\mathfrak{u}(n)$ 的基本向量场，在令 $g(t)=\exp(tY)$ 时为
  \begin{align}
    \underline{Y}&=\sum_{i=1}^n\left.\frac{d(g(t)\cdot z)_i}{dt}\right|_{t=0}\frac{\partial}{\partial z_i}+(\text{c.c.}) \tag{5.1.52}\\
    &=\sum_{i=1}^n\left.\frac{d(e^{tY}\cdot z)_i}{dt}\right|_{t=0}\frac{\partial}{\partial z_i}+(\text{c.c.}) \tag{5.1.53}\\
    &=\sum_{i=1}^n\left.(Ye^{tY}\cdot z)_i\right|_{t=0}\frac{\partial}{\partial z_i}+(\text{c.c.}) \tag{5.1.54}\\
    &=\sum_{i=1}^n\left(\sum_{j=1}^n Y_{ij}z_j\right)\frac{\partial}{\partial z_i}+(\text{c.c.}) \tag{5.1.55}\\
    &=\sum_{i,j=1}^n Y_{ij}z_j\frac{\partial}{\partial z_i}+\sum_{i,j=1}^n\overline{Y_{ij}}\,\bar z_j\frac{\partial}{\partial\bar z_i}. \tag{5.1.56}
  \end{align}
  其中，(c.c.) 表示复共轭（complex conjugate）。因此，
  % [原文 p.64]
  \begin{align}
    \mu_Y&=-\iota_{\underline{Y}}\theta \tag{5.1.57}\\
    &=-\frac{\sqrt{-1}}{4}\sum_{i,j=1}^n(Y_{ij}z_j\bar z_i-\overline{Y_{ij}}\,\bar z_jz_i) \tag{5.1.58}\\
    &=-\frac{\sqrt{-1}}{2}\sum_{i,j=1}^n Y_{ij}z_j\bar z_i, \tag{5.1.59}
  \end{align}
  其中最后的等号用到由
  \begin{equation}
    \mathfrak{u}(n)=\{Y\in\mathrm{Mat}_{n\times n}(\C)\mid Y^*=-Y\}
    \tag{5.1.60}
  \end{equation}
  得到的
  \begin{equation}
    \overline{Y_{ij}}\,\bar z_jz_i=-Y_{ji}z_i\bar z_j
    \tag{5.1.61}
  \end{equation}
  以及
  \begin{equation}
    \sum_{i,j=1}^n Y_{ji}z_i\bar z_j=\sum_{i,j=1}^n Y_{ij}z_j\bar z_i.
    \tag{5.1.62}
  \end{equation}
  此处，用内积
  \begin{align}
    (Y,Z)&:=-\operatorname{tr}(YZ) \tag{5.1.63}\\
    &=-\sum_{i,j=1}^n Y_{ij}Z_{ji} \tag{5.1.64}\\
    &=\sum_{i,j=1}^n Y_{ij}\overline{Z_{ij}} \tag{5.1.65}
  \end{align}
  把 $\mathfrak{g}$ 与 $\mathfrak{g}^*$ 视为同一，则
  \begin{equation}
    \mu_Y=-\operatorname{tr}(Y\mu)
    \tag{5.1.66}
  \end{equation}
  由此得到
  \begin{equation}
    \mu(z)=\left(\frac{\sqrt{-1}}{2}z_i\bar z_j\right)_{i,j=1}^n.
    \tag{5.1.67}
  \end{equation}
\end{proof}
% [Tutor Note: 难点——(5.1.54) 是把 $d(e^{tY}\cdot z)/dt=(Ye^{tY})\cdot z$ 写成矩阵乘法后取 $t=0$；(5.1.61) 来自 $Y^*=-Y$ 即 $\overline{Y_{ij}}=-Y_{ji}$。]

关于动量映射，下述定理是基本的：

\begin{theorem}{5.1.9（[10], [55], [57]）}
  设环面 $T^n=(S^1)^n$ 对紧致辛流形 $(M,\omega)$ 有 Hamilton 作用，其动量映射由
  \begin{equation}
    \mu\colon M\to\mathfrak{t}^*\cong\R^n
    \tag{5.1.68}
  \end{equation}
  给出。此时，动量映射的任意纤维都是连通的，其像是以 $T^n$ 作用的固定点的凸包形式给出的凸多面体。
\end{theorem}

证明用到 Morse 理论，此处省略。

% [Tutor Note: 背景——此即 Atiyah 与 Guillemin--Sternberg 的凸性定理（1982）；[10]、[55]、[57] 为相关原始文献，证明须用 Morse 理论，故原著从略。]

% [原文 p.65]
\section{辛商}

给出紧致 Lie 群对辛流形的 Hamilton 作用后，可按如下方式定义与之相伴的商：

\begin{theorem*}{5.2.1}
  设紧致 Lie 群 $G$ 对辛流形 $(M,\omega)$ 有 Hamilton 作用，其动量映射由 $\mu\colon M\to\mathfrak{g}^*$ 给出。进一步设 $0\in\mathfrak{g}^*$ 是 $\mu$ 的正则值，且 $G$ 对 $\mu^{-1}(0)$ 的作用自由；
  \begin{equation}
    \begin{array}{ccc}
      \mu^{-1}(0) & \xrightarrow{\ \iota\ } & M\\[2pt]
      {\scriptstyle\pi}\big\downarrow & & \\[2pt]
      \mu^{-1}(0)/G & &
    \end{array}
    \tag{5.2.1}
  \end{equation}
  此时，$M/\!\!/G:=\mu^{-1}(0)/G$ 上满足
  \begin{equation}
    \pi^*\omega_{\mathrm{red}}=\iota^*\omega_M
    \tag{5.2.2}
  \end{equation}
  的辛形式 $\omega_{\mathrm{red}}$ 唯一存在。
\end{theorem*}

\begin{proof}
  因 $0\in\mathfrak{g}^*$ 是 $\mu$ 的正则值，$\mu^{-1}(0)$ 是 $M$ 的子流形。因 $G$ 对 $\mu^{-1}(0)$ 的作用自由，$M/\!\!/G:=\mu^{-1}(0)/G$ 是流形。把 (5.2.2) 换个说法，即：对任意 $p\in\mu^{-1}(0)$ 与任意 $X,Y\in T_p\mu^{-1}(0)$ 有
  \begin{equation}
    (\omega_{\mathrm{red}})_{\pi(p)}(\pi_*X,\pi_*Y)=\omega_M(\iota_*X,\iota_*Y).
    \tag{5.2.3}
  \end{equation}
  由于 $\pi$ 是满射，满足此式的 $M/\!\!/G$ 上的 2 次微分形式只要存在便唯一。为证存在，对任意 $q\in M/\!\!/G$ 与 $\bar X,\bar Y\in T_q(M/\!\!/G)$，取满足 $\pi(p)=q$ 的 $p\in\mu^{-1}(0)$ 与满足 $\pi_*X=\bar X$、$\pi_*Y=\bar Y$ 的 $X,Y$，按
  \begin{equation}
    (\omega_{\mathrm{red}})_q(\bar X,\bar Y)=(\omega_M)_p(\iota_*X,\iota_*Y)
    \tag{5.2.4}
  \end{equation}
  % [Tutor Note: 校勘——原书写「不依赖于 $q\in\pi^{-1}(p)$ 的取法」，按文意应为 $p\in\pi^{-1}(q)$（下文正是按 $p$ 在纤维中的选取来论证）；此处照录。]
  定义后，须证它是良定义的（即不依赖于 $q\in\pi^{-1}(p)$ 与 $X,Y\in T_p(\mu^{-1}(0))$ 的取法）。换言之，须证：对任意 $p,p'\in\mu^{-1}(0)$ 与 $X,Y\in T_p(\mu^{-1}(0))$、$X',Y'\in T_{p'}(\mu^{-1}(0))$，当 $\pi(p)=\pi(p')$、$\pi_*X=\pi_*X'$ 且 $\pi_*Y=\pi_*Y'$ 时，$\omega_M(X,Y)=\omega_M(X',Y')$。首先，由 $\pi(p)=\pi(p')$，存在 $g\in G$ 使 $p=g\cdot p'$。用它令 $X''=g_*X'$、$Y''=g_*Y'$，则
  \begin{equation}
    \pi_*X'=(\pi\circ g)_*X'=\pi_*(g_*X')=\pi_*X'',\qquad \pi_*Y'=\pi_*Y''
    \tag{5.2.5}
  \end{equation}
  成立。又，因 $\omega$ 在 $G$ 作用下不变，
  % [原文 p.66]
  \begin{equation}
    (\omega_M)_{p'}(X',Y')=(\omega_M)_p(X'',Y'')
    \tag{5.2.6}
  \end{equation}
  成立。这里，$X,Y$ 与 $X'',Y''$ 都是 $p$ 处 $\pi^{-1}(0)$ 的切向量，且由 $\pi_*X=\pi_*X''$ 得 $X-X''\in\Ker\pi_*$。由于 $\pi\colon\mu^{-1}(0)\to\mu^{-1}(0)/G$ 是自由 $G$ 作用下的商，这意味着存在某个 $Z\in\mathfrak{g}$，使
  \begin{equation}
    X-X''=\underline{Z}_p
    \tag{5.2.7}
  \end{equation}
  成立。此时，对任意 $W\in T_p(\mu^{-1}(0))$ 有
  \begin{align}
    \omega(\underline{Z}_p,W)&=d\mu_Z(W) \tag{5.2.8}\\
    &=(d\langle\mu,Z\rangle)(W) \tag{5.2.9}\\
    &=\langle d\mu(W),Z\rangle \tag{5.2.10}\\
    &=0, \tag{5.2.11}
  \end{align}
  其中最后的等号用到由 $W\in T_p(\mu^{-1}(0))$ 得出的
  \begin{equation}
    d\mu(W)=0
    \tag{5.2.12}
  \end{equation}
  于是，
  \begin{equation}
    (\omega_M)_p(X'',Y'')=(\omega_M)_p(X,Y'')
    \tag{5.2.13}
  \end{equation}
  同样可得
  \begin{equation}
    (\omega_M)_p(X,Y'')=(\omega_M)_p(X,Y),
    \tag{5.2.14}
  \end{equation}
  从而推出 $\omega_{\mathrm{red}}$ 是良定义的。

  由 $\omega_M$ 的非退化性与
  \begin{equation}
    \dim\mu^{-1}(0)=\dim M-\dim G
    \tag{5.2.15}
  \end{equation}
  可知，$\omega|_{\mu^{-1}(0)}$ 的核的维数至多为 $\dim G$。另一方面，$G$ 作用的基本向量场张成 $T(\mu^{-1}(0))$ 的一个子丛，其维数为 $\dim G$，并且包含于 $\omega|_{\mu^{-1}(0)}$ 的核中。因此，$\omega|_{\mu^{-1}(0)}$ 的核由 $G$ 作用的基本向量场张成，从而可知 $\omega_{\mathrm{red}}$ 是非退化的。
\end{proof}
% [Tutor Note: 校勘——(5.2.6) 之后原书把 $X,Y,X'',Y''$ 说成「$p$ 处 $\pi^{-1}(0)$ 的切向量」，但 $\pi$ 的定义域是 $\mu^{-1}(0)$（下文即用 $\pi$ 的商映射与 $\Ker\pi_*$），故 $\pi^{-1}(0)$ 应为 $\mu^{-1}(0)$，此处照录。]
% [Tutor Note: 难点——非退化性：$\omega$ 限制到 $\mu^{-1}(0)$ 后核的维数至多为 $\dim G$（由 (5.2.15)），而基本向量场张成的 $\dim G$ 维子丛已含于核中，故核恰为该子丛，商掉的正是它。]

\takeaway{在动量取正则值处按群作用取商，辛形式沿轨道收缩，即得辛商 $M/\!\!/G=\mu^{-1}(0)/G$。}

更一般地，对余伴随轨道 $\mathcal{O}\subset\mathfrak{g}^*$，若考虑以把 Kostant--Kirillov 形式的符号取反所得者为辛形式的辛结构，则动量映射成为给包含映射加上负号者。它与辛流形 $(M,\omega)$ 的直积 $(M\times\mathcal{O},\omega_M-\omega_{\mathcal{O}})$ 具有对角型的 Hamilton 作用
\begin{equation}
  G\curvearrowright M\times\mathcal{O},\qquad g\colon(m,\xi)\mapsto(g\cdot m,\operatorname{Ad}^*_g\xi),
  \tag{5.2.16}
\end{equation}
% [原文 p.67]
其动量映射由
\begin{equation}
  \mu-\iota\colon(m,\xi)\mapsto\mu(m)-\xi
  \tag{5.2.17}
\end{equation}
给出。此时记
\begin{align}
  M/\!\!/_{\mathcal{O}}G&:=(M\times\mathcal{O})/\!\!/G \tag{5.2.18}\\
  &=\{(m,\xi)\in M\times\mathfrak{g}^*\mid\mu(m)-\xi=0\}/G \tag{5.2.19}\\
  &=\mu^{-1}(\mathcal{O})/G, \tag{5.2.20}
\end{align}
并称之为 $M$ 关于 $\mathcal{O}\subset\mathfrak{g}^*$ 的\emph{辛商}（symplectic quotient），或 \emph{Marsden--Weinstein 约化}（Marsden--Weinstein reduction）。替换 $\mathcal{O}$，便可对同一 Hamilton 作用考虑种种商。理解它们之间的关系，在理论上与应用上都是既有意义又重要的问题。
% [Tutor Note: 背景——Marsden--Weinstein 约化由 Marsden 与 Weinstein 于 1974 年提出，Meyer 亦于 1973 年独立得到；对余伴随轨道取商把「动量取固定值」推广为「动量落在轨道内」。]

\chapterreview{本章以 Noether 定理为线索，说明辛几何中正确的商是「先固定动量、再按群作用取商」的辛商。5.1 节定义 Hamilton 作用与动量映射，用等变性刻画其 Poisson 映射条件，并给出闭子群、直积、余伴随轨道、恰当辛流形、余切丛与 $U(n)$ 线性作用等典型例子，最后陈述环面作用的凸性定理。5.2 节证明 Marsden--Weinstein 约化定理：在 $0$ 为正则值、作用自由的假设下，$M/\!\!/G=\mu^{-1}(0)/G$ 上存在唯一的辛形式满足 $\pi^*\omega_{\mathrm{red}}=\iota^*\omega_M$。更一般地，对任意余伴随轨道 $\mathcal{O}$ 可作辛商 $M/\!\!/_{\mathcal{O}}G=\mu^{-1}(\mathcal{O})/G$，它与 GIT 商、VGIT 的联系是后续诸章（尤其第 8 章的环面几何）反复出现的主题。}

